\par\Needspace{4\baselineskip}
\begin{center}
\SourceLeaf{040}\pdfbookmark[2]{22.}{article-022}22.
\end{center}
Differentiating formula IV gives (with $\psi$ treated as constant)
\[
\begin{aligned}
\frac{\mathrm{d}u}{u}&=\frac{1}{2}\left(\operatorname{tang}\frac{1}{2}(v+\psi)-\operatorname{tang}\frac{1}{2}(v-\psi)\right)\mathrm{d}v\\
&=\frac{r\operatorname{tang}\psi}{p}\,\mathrm{d}v;
\end{aligned}
\]
hence
\[
r r\,\mathrm{d}v=\frac{p r}{u\operatorname{tang}\psi}\,\mathrm{d}u
\]
or, substituting the value of $r$ from X,
\[
r r\,\mathrm{d}v=b b\operatorname{tang}\psi\left\{\frac{1}{2}e\left(1+\frac{1}{u u}\right)-\frac{1}{u}\right\}\mathrm{d}u.
\]
Integrating so that the integral vanishes at perihelion then gives
\[
\begin{aligned}
\int r r\,\mathrm{d}v&=b b\operatorname{tang}\psi\left\{\frac{1}{2}e\left(u-\frac{1}{u}\right)-\log u\right\}\\
&=k t\BareRoot{p}\, .\,\BareRoot{(1+\mu)}=k t\operatorname{tang}\psi\, .\,\BareRoot{b}\, .\,\BareRoot{(1+\mu)}.
\end{aligned}
\]
The logarithm here is hyperbolic. If we choose to use logarithms from the \textsc{Briggs}ian system or, more generally, from a system whose modulus is $=\lambda$, and if we neglect the mass $\mu$ (which we may suppose cannot be determined for a body moving in a hyperbola), the equation takes the following form.
\[
\mathrm{XI.}\quad\begin{gathered}
\frac{1}{2}\lambda e\frac{u u-1}{u}-\log u=\frac{\lambda k t}{b^{\frac{3}{2}}},\\
\text{or, introducing }F\\
\lambda e\operatorname{tang}F-\log\operatorname{tang}(45^0+\frac{1}{2}F)=\frac{\lambda k t}{b^{\frac{3}{2}}}\cdot
\end{gathered}
\]
If we assume that \textsc{Briggs}ian logarithms are used, we have \\
$\log\lambda=9{,}637\,7843\,113$, $\log\lambda k=7{,}873\,3657\,527$, but slightly greater precision can be attained by using hyperbolic logarithms directly. Hyperbolic logarithms of tangents are found in several collections of tables, for example in the one prepared by \textsc{Schulze}, and over a still greater range in the Magnus Canon Triangulorum Logarithmicus of \textsc{Beni. Ursinus}, Colon. 1624, where they are tabulated at intervals of $10^{\prime\prime}$. — Formula XI also shows that reciprocal values of $u$, or opposite values of $F$ and $v$, correspond to opposite values of $t$. Thus equal portions of the hyperbola at equal distances on either side of perihelion are traversed in equal times.
